\label{sec:presaAria}
\section{Diffusore subsonico}
Il diffusore subsonico è stato modellato come un componente adiabatico,
assumendo condizioni uniformi nelle sezioni caratteristiche e imponendo
$M_2=0.5$ all'ingresso del compressore. Il recupero di pressione totale è stato determinato mediante il rendimento isentropico del diffusore, ricavando il corrispondente coefficiente $\pi_D$. L'area di cattura è stata assunta coincidente con quella geometrica di ingresso ($A_0=A_1$).

Il flusso destinato al turbogetto è stato separato da quello del ramjet mediante un setto mobile interno, mantenendo una presa geometrica comune. Note le aree di ingresso e uscita ottenute dall'ottimizzazione, il diffusore è stato rappresentato mediante una geometria conica a sezione circolare. È stato assunto il rapporto minimo $L/D_1=1.5$ al fine di contenere la lunghezza del componente; la geometria risultante è stata successivamente verificata rispetto ai limiti di separazione del flusso. Il dimensionamento e la relativa verifica sono riportati in Appendice~\ref{app:presa_aria_subsonica}.


% \usepackage{tikz}
% \usetikzlibrary{arrows.meta,angles,quotes,calc}
% \usepackage{booktabs}

\begin{figure}[H]
\centering

%========================================================
% DIFFUSORE A SINISTRA
%========================================================
\begin{minipage}[c]{0.72\textwidth}
\centering

\begin{tikzpicture}[
    x=5.2cm,
    y=5.2cm,
    >=Latex,
    line cap=round,
    line join=round
]

%------------------------------------------------
% DATI GEOMETRICI
%------------------------------------------------
\def\L{1.105}
\def\Dzero{0.737}
\def\Dtwo{0.866}

\pgfmathsetmacro{\Rzero}{\Dzero/2}
\pgfmathsetmacro{\Rtwo}{\Dtwo/2}

%------------------------------------------------
% COORDINATE PRINCIPALI
%------------------------------------------------
\coordinate (A) at (0,\Rzero);
\coordinate (B) at (0,-\Rzero);

\coordinate (C) at (\L,\Rtwo);
\coordinate (D) at (\L,-\Rtwo);

%------------------------------------------------
% PROFILO DEL DIFFUSORE
%------------------------------------------------
\draw[line width=1.1pt] (A) -- (C);
\draw[line width=1.1pt] (B) -- (D);

\draw[line width=1.1pt] (A) -- (B);
\draw[line width=1.1pt] (C) -- (D);

% Asse centrale
\draw[dash dot, line width=0.6pt]
(-0.05,0) -- (\L+0.05,0);

% Linee di sezione
\draw[dash dot, line width=0.6pt]
(0,-0.52) -- (0,0.58);

\draw[dash dot, line width=0.6pt]
(\L,-0.52) -- (\L,0.58);

%------------------------------------------------
% ETICHETTE SEZIONI
%------------------------------------------------
\node[align=center] at (0,0.62)
{Sezione 0\\(ingresso)};

\node[align=center] at (\L,0.62)
{Sezione 2\\(uscita)};

%------------------------------------------------
% FRECCIA DEL FLUSSO
%------------------------------------------------
\draw[->, line width=0.9pt]
(0.25,0.54) -- (0.45,0.54);

\node at (0.35,0.59) {Flusso};

%------------------------------------------------
% QUOTA D0
%------------------------------------------------
\draw[line width=0.5pt]
(-0.08,\Rzero) -- (-0.01,\Rzero);

\draw[line width=0.5pt]
(-0.08,-\Rzero) -- (-0.01,-\Rzero);

\draw[<->, line width=0.6pt]
(-0.07,-\Rzero) -- (-0.07,\Rzero);

\node[left] at (-0.07,0)
{$D_0=0.737\,\mathrm{m}$};

%------------------------------------------------
% QUOTA D2
%------------------------------------------------
\draw[line width=0.5pt]
(\L+0.01,\Rtwo) -- (\L+0.08,\Rtwo);

\draw[line width=0.5pt]
(\L+0.01,-\Rtwo) -- (\L+0.08,-\Rtwo);

\draw[<->, line width=0.6pt]
(\L+0.07,-\Rtwo) -- (\L+0.07,\Rtwo);

\node[right] at (\L+0.07,0)
{$D_2=0.866\,\mathrm{m}$};

%------------------------------------------------
% QUOTA L
%------------------------------------------------
\draw[line width=0.5pt]
(0,-0.49) -- (0,-0.43);

\draw[line width=0.5pt]
(\L,-0.49) -- (\L,-0.43);

\draw[<->, line width=0.6pt]
(0,-0.47) -- (\L,-0.47);

\node[fill=white, inner sep=2pt]
at (\L/2,-0.47)
{$L=1.105\,\mathrm{m}$};

%------------------------------------------------
% ANGOLO PHI
%------------------------------------------------
\draw[dashed, line width=0.5pt]
(0.05,\Rzero) -- (0.32,\Rzero);

\draw[->, line width=0.6pt]
(0.23,\Rzero+0.005)
arc[start angle=0,end angle=3.35,radius=0.18];

\node at (0.33,\Rzero+0.055)
{$\phi=3.35^\circ$};

\end{tikzpicture}

\end{minipage}
\hfill
%========================================================
% TABELLA A DESTRA
%========================================================
\begin{minipage}[c]{0.25\textwidth}
\centering

\footnotesize
\renewcommand{\arraystretch}{0.90}
\setlength{\tabcolsep}{3pt}

\begin{tabular}{lcc}
\toprule
\textbf{Parametro} & \textbf{Valore} & \textbf{Unità} \\
\midrule

$\pi_D$ & 0.989  & -- \\
$A_0$   & 0.4262 & $\mathrm{m^2}$ \\
$A_2$   & 0.5889 & $\mathrm{m^2}$ \\

\bottomrule
\end{tabular}

\end{minipage}

\caption{Geometria e principali parametri del diffusore subsonico nel punto di progetto.}
\label{fig:geometriaDiffusore}

\end{figure}